\documentclass[a4paper]{article}     
\usepackage{amsfonts}
\usepackage{amsmath}
\usepackage{amssymb}
\usepackage{graphicx}


\begin{document}

\noindent \large Auteur: Xander De Blauwer \\
\noindent \large Studierichting: Informatica\\
\noindent \large Oefeningenreeks 4A nr. 63.5\\

\medskip

\normalsize


$ = \displaystyle\int \dfrac{dx}{(\sqrt{1+x^2})^3}$ \\

Stel $x = \tan(t) \Rightarrow dx = \sec^2(t)dt $ \\

$ = \displaystyle\int \dfrac{\sec^2(t)}{\sqrt{\sec^2(t)}^3}dt $ \\

$ = \displaystyle\int \dfrac{dt}{\sec(t)} = \displaystyle\int \cos(t)dt $ \\

$ = \sin(t) + c = \sin(\operatorname{Bgtan}(x)) + c $ \\

Stel $\theta = \operatorname{Bgtan}(x) \Rightarrow \tan(\theta) = x $ \\

$\tan(\theta) = \dfrac{\sin(\theta)}{\cos(\theta)} $ \\

$ \Rightarrow x = \dfrac{\sin(\theta)}{\cos(\theta)} \Leftrightarrow \sin(\theta) = x\cos(\theta) $ \\

$x^2\cos^2(\theta) + \cos^2(\theta) = 1 $ \\

$\Rightarrow \cos(\theta) = \dfrac{1}{\sqrt{1+x^2}}$ \\

$\Rightarrow \sin(\theta) = \dfrac{x}{\sqrt{1+x^2}}$

$\displaystyle\int \dfrac{dx}{(\sqrt{1+x^2})^3} = \dfrac{x}{\sqrt{1+x^2}} + c $ \\


\begin{thebibliography}{999}
%\bibitem{a} Referentie
\end{thebibliography}
\end{document} 